% Study 3: supplementary methods and the checks and text analyses that Supplementary Note 3 does
% not give. Numbers come only from macros in generated/numbers.tex.
\label{sec:suppB-study3}

Study~3 asked whether stating how long Alex spent on an original closes the gap between the ratings given under the two labels. Its data were collected from 4 to 9 March 2025. Its registration, a dated record of its hypotheses and analysis plan, was filed on the Open Science Framework (OSF) on 8 October 2026, after an April 2025 analysis of the same data and a September 2025 talk (Supplementary Table~\ref{tab:supp-verdicts}). Supplementary Note~3 gives every registered test, the gap at each stated time and the predicted means. This section gives the materials, procedure and sample, and the checks and text analyses that Note~3 does not.

\subsection*{Materials}

Study~3 used Study~1's ten originals. All ten were generated with OpenAI's DALL·E, an AI image generator (Section~\ref{sec:suppA-registrations}), so the Human label was a cover story, and five show a human or humanoid figure. Each participant saw nine of the ten, drawn at random, and each was shown alone for at least 8 seconds. There were no second works.

The Human label read: `There is a digital artist named Alex. Alex is a skilled artist who creates digital artworks. Alex uses only his graphic design skills and digital sketching tools such as iPad apps and Adobe Photoshop.' The AI label, the same in all five studies, read: `There is a digital artist named Alex. Alex is a skilled artist who creates digital artworks. Alex develops and inputs prompts into an artificial intelligence (AI) system. The AI system generates digital artworks based on those prompts.'

Before each original, the survey inserted a stated time, `less than 1 hour', `10 hours' or `100 hours', where \textit{time} appears below. The Human group read `Alex had a new idea for a digital artwork. Alex spent \textit{time} working only with his graphic design and digital sketching tools, revising as his work as he goes, which creates a digital artwork.' The AI group read `Alex had a new idea for a work of digital artwork. Alex spent \textit{time} developing and inputting prompts into the AI image generator, revising each prompt as he goes, which creates a digital artwork.' Both went on: `Click the arrow to see the digital artwork created. Note that you will be required to stay on the page for at least 8 seconds before proceeding.' The quotations reproduce the survey exactly, slips included.

Two features of this wording belong to Study~3 alone. Its Human label is the only one that names the tools. And of the two studies that asked about originals, it is the only one whose vignettes state a time. Under the Human label, the hours were spent drawing; under the AI label, they were spent revising prompts, so there more time also meant more control over the image, and the two cannot be told apart.

On the screen after each original, which no longer showed the image, participants rated creativity and then effort, on five-point scales scored 1 to 5 by the position of the answer chosen. The creativity item read `How much creativity was involved in the production of this work? Creativity is defined as the presence of at least a minimal original, intellectual contribution to the creation of the work.' Its answers ran from `No creativity involved' through minimal, moderate and significant creativity to `A great deal of creativity involved'. The effort item read `How much effort was involved in the production of this work?', with no definition, and its answers ran from `No effort involved' through minimal, moderate and significant effort to `A great deal of effort involved'. Their wording matches Study~1's, except that one of Study~1's effort items had different answer labels; in Study~3 the effort answers were the same for every image. Below the two ratings, a required text box for each asked `Please explain your reasoning for your rating on creativity:' (or `on effort:'). Each explanation was therefore written after its rating had been given, although a rating could still be changed before the participant left the screen. No box blocked copying or pasting.

\subsection*{Procedure}

Participants took part online in a Qualtrics survey. By our account, participants were students in the UCLA Psychology subject pool who took part for course credit. Every analyzed record carries the code that the survey sent, at its end, to the pool's page for granting credit. The consent form, the same in all five studies, also offered hourly pay to people who signed up outside the pool, and no record shows that anyone took that route. The survey assigned each participant at random to one label for the whole session, presenting each label equally often. It also set a quota of 50 participants for each label group, a planned number that neither stopped anyone nor moved anyone to the other label. A script attached a stated time to each trial, drawing at random among the times not yet used three times, so each participant saw three originals at each time in a random order. On each trial, participants read the vignette with its stated time and saw the original for at least 8 seconds; on the next screen they gave both ratings and then wrote both explanations: 18 ratings and 18 explanations in all.

A question about the task (`What is your task in this experiment?') came before the label and allowed one retry, so it ended no session. An arithmetic question (`What does 8+7 equal?') came after the trials. A closing screen then asked first whether the participant had completed the experiment seriously, without responding randomly, then whether any part of the procedure had been unclear, and then about race, gender, age, native language and formal artistic training. No question asked participants to recall the label (a manipulation check) or whether they believed it. The debriefing was the same as in the other four studies; Section~\ref{sec:suppA-study1} quotes it and corrects its statement of the law. Participants read it only after their last rating.

\subsection*{Exclusions and sample}

The registered exclusion rules, applied in order, left \SThreeN{} of the \SThreeNRecords{} records in the survey export: \SThreeNPreviewRemoved{} survey test runs were removed, then \SThreeNUnfinishedRemoved{} unfinished records, \SThreeNNotSeriousRemoved{} participants who gave any answer but `Yes, I did' to the seriousness question, a blank included, and \SThreeNArithmeticRemoved{} who answered the arithmetic question wrongly. Of the \SThreeN{} participants analyzed, \SThreeNHuman{} saw the Human label and \SThreeNAI{} the AI label. Of their \SThreeNJustifications{} explanations, \SThreeNRatings{} came with a scored rating. The few remaining ratings were blank or carried an answer code that the final survey did not offer, and they were treated as missing.

The participants' mean age was \SThreeAgeMean{} years (standard deviation \SThreeAgeSD; range \SThreeAgeMin{} to \SThreeAgeMax). Of the \SThreeN, \SThreeNWomen{} described their gender as female, \SThreeNMen{} as male, and \SThreeNGenderOther{} chose `Other' or `Prefer not to say' or gave no answer. The answers on race, native language and formal artistic training are not summarized here.

\subsection*{Results not reported elsewhere}

Each result below carries its status. A registered confirmatory test is one that the registration names as confirmatory, judged by its rule after Holm's correction, which raises each $p$-value in a set of tests so that the chance of calling any of them significant by mistake stays at .05. A $p$-value is the chance of a result at least this large if there were no effect, and a result is called significant when $p$ is below .05. A registered secondary analysis is one the registration names outside that set. An exploratory analysis was named as exploratory or not named at all, and carries no verdict.

\paragraph*{The ordinal check.}
The registration named a cumulative-link mixed model as a check on each confirmatory verdict (registered secondary). This model treats the five answers as ordered categories rather than evenly spaced numbers, and it tests each term with a chi-square statistic, corrected by Holm's method over the same six tests as the linear model. On creativity, the label gave $\chi^2$~=~\SThreeCreatSourceOrdChisq{} (Holm $p$~\SThreeCreatSourceOrdPHolm), stated time \SThreeCreatTimeOrdChisq{} (Holm $p$~\SThreeCreatTimeOrdPHolm), and their interaction \SThreeCreatSourceByTimeOrdChisq{} (Holm $p$~\SThreeCreatSourceByTimeOrdPHolm). On effort, the label gave \SThreeEffSourceOrdChisq{} (Holm $p$~\SThreeEffSourceOrdPHolm) and stated time \SThreeEffTimeOrdChisq{} (Holm $p$~\SThreeEffTimeOrdPHolm); Supplementary Note~3 gives the effort interaction, which came nearest to the threshold. Every verdict matched the linear model's. Both registered linear models and both ordinal models converged, meaning that the fitting procedure settled on a solution, and none of the four was singular, that is, none estimated a source of variation at zero. The registered refit that also lets the time effect vary from person to person and the label's effect vary from image to image was singular on both ratings, so, as registered, its requirement was waived (Supplementary Note~3).

\paragraph*{The contrast that set Study~5's bound.}
Averaged over labels, the creativity mean at 100 hours minus the mean at under 1 hour was \SThreeTimeContrastPrecise{} points before rounding (registered secondary; Supplementary Note~3). Before fitting any of its own models, Study~5 took this value as its bound for stated editing time, $\pm$\SFiveEditTimeBound{} points.

\paragraph*{Registered text checks.}
Two checks on the explanations were registered for every study, and Study~3 ran both on creativity and on effort (registered secondary; Section~\ref{sec:suppA-text} describes both). In the first, ridge regression, a regression that shrinks each word's weight toward zero, predicted each rating and each participant's label from the words of explanations written by participants left out of the fit. The words predicted the rating: predicted and actual ratings correlated at \SThreeTextRecRatingCreat{} on creativity (95\% confidence interval, CI, the range of values the data do not rule out at the 95\% level: \SThreeTextRecRatingCreatLo{} to \SThreeTextRecRatingCreatHi) and \SThreeTextRecRatingEff{} on effort (\SThreeTextRecRatingEffLo{} to \SThreeTextRecRatingEffHi). In runs with the ratings shuffled, 95\% of correlations stayed below \SThreeTextRecRatingCreatNull{} and \SThreeTextRecRatingEffNull{} respectively. The label was shown by the effort explanations but not by the creativity explanations. The AUC, the chance that a randomly chosen AI-label explanation scores above a randomly chosen Human-label one, so that 0.5 is chance, was \SThreeTextRecArmEff{} on effort (\SThreeTextRecArmEffLo{} to \SThreeTextRecArmEffHi) and \SThreeTextRecArmCreat{} on creativity (\SThreeTextRecArmCreatLo{} to \SThreeTextRecArmCreatHi, an interval that reaches chance). The rating alone gave \SThreeTextRecArmFromRatingEff{} and \SThreeTextRecArmFromRatingCreat. Adding the words to the rating changed the AUC by \SThreeTextRecIncrEff{} on effort (\SThreeTextRecIncrEffLo{} to \SThreeTextRecIncrEffHi; the lower end lies just below zero before rounding) and \SThreeTextRecIncrCreat{} on creativity (\SThreeTextRecIncrCreatLo{} to \SThreeTextRecIncrCreatHi), so on both the registered reading is inconclusive. The second check predicted that AI-label participants would reuse their own words across explanations more than Human-label participants, measured as a person's word overlap among their own explanations minus their overlap with other people's explanations of the same image, with explanation length held constant, and Holm's correction ran across the two ratings. The difference, AI minus Human, was \SThreeTextOverlapCreatDiff{} on creativity (\SThreeTextOverlapCreatLo{} to \SThreeTextOverlapCreatHi; Holm $p$~\SThreeTextOverlapCreatPHolm) and \SThreeTextOverlapEffDiff{} on effort (\SThreeTextOverlapEffLo{} to \SThreeTextOverlapEffHi; Holm $p$~\SThreeTextOverlapEffPHolm), so the prediction is inconclusive on both, which is not evidence that the groups reuse words equally. Every explanation was written after its rating, so these checks show what the words track, not why anyone rated as they did.

\paragraph*{Exploratory text analyses.}
Supplementary Note~7 gives the word-list comparisons between labels. Compared across stated times instead, without correction, effort explanations at 100 hours mentioned time more often, and prompting and machines less often, than at under 1 hour, and creativity explanations at 100 hours mentioned visible detail more often. Participants often restated the stated time, so a mention of time may echo the vignette. Word tone and concreteness, scored from published norms, did not differ by label after a Benjamini--Hochberg correction, which limits the expected share of false positives among results called significant. Two topics found without labels, groups of words that tend to occur together, leaned toward the AI label before correction: a creativity topic led by `ai', `created' and `generated', and an effort topic led by `hours', `spent' and `prompt'. Neither survived correction. The words that carried the prediction of the rating mostly echoed the answer labels, such as `minimal' and `great'. Their estimates are in the Study~3 tables that the analysis code writes (main text, Code availability).
